\cotefile{134-2}
\scanpage{170}\mspage{25}%
% archive sheet number: 170
On dira qu'une $n$-catégorie \inserted{$\unsure{\mathcal E|E}_{n}$} est ``\uline{spéciale}'' (ou $n$-\uline{galoisienne}) si elle est $n$-équivalente à la cat. des $(n-1)$-champs loc. constants sur un espace top. (ou un topos) convenable, ou ce qui doit revenir au $\unsure{=|\overline m}$, si elle est $n$-équivalente à la catégorie des $n$-foncteurs $G_n\longrightarrow(n\text{-}\mathrm{Cat})$, où $G_n$ est un $n$-groupoïde, Si $X$, $G_n$, $\mathcal E_n$ se correspondent ainsi, on dit que $\deleted{\unsure{G_n|\mathcal E_n}}\,G_n$ est le $n$-groupoïde fondamental de $X$, ou de $\mathcal E_n$, ou que $\mathcal E_n$ est la catégorie des $n$-systèmes locaux sur $X$, ou sur $G_n$, ou que $X$ est une réalisation géométrique \unsure{de|du} $G_n$ ou de $\mathcal E_n$. En analogie avec le cas bien familier $n=1$, $G_n$ doit s'interpréter comme la sous-$n$-catégorie pleine de $\underline{\mathrm{Hom}}_n(\mathcal E_n,(n\text{-}\mathrm{cat}))$ formée des $n$-foncteurs $\mathcal E_n\to(n\text{-}\mathrm{cat})$ satisfaisant certaines propriétés d'exactitude (on a envie de dire : qui commutent aux $\unsure{\varprojlim|\underline{\lim}}$ finies et aux $\varinjlim$ quelconques) ; mais cela nous
% Faint interlinear curved mark below nous, whole-image x290–318, y1991–2018; hand/function uncertain.
\raisebox{-1.8ex}[0pt][0pt]{\makebox[0pt][r]{\parbox{3mm}{\figure{\begin{tikzpicture}[x=1mm,y=1mm]
\draw[line width=0.25pt,gray] (0,0) .. controls (0.35,-1.15) and (1.55,-2.45) .. (2.2,-2.55) .. controls (2.65,-2.7) and (2.7,-2.2) .. (2.38,-1.85);
\end{tikzpicture}}}}}
 attire dans l'inquiétant engrenage des $n$-limites dans les $n$-catégories (\uline{NB} le cas $n=2$ commence à nous être familier\ldots) Il est prudent dans tout
